% Auto-generated by generate_assets.py --- tab_bis.tex
\begin{table}[htbp]
  \centering
  \small
  \caption{Score d'impact du biais (BIS), calculé comme $\mathrm{BIS}=\mathrm{CR}-\mathrm{RR}$ avec le CR du même biais. Un BIS positif indique que le juge change plus souvent de verdict après perturbation qu'entre deux exécutions identiques.}
  \label{tab-bis}
  \sisetup{table-number-alignment=center}
  \begin{tabular}{lSSS}
    \toprule
    \textbf{Biais} & {\textbf{DeepSeek}} & {\textbf{Gemma}} & {\textbf{Qwen3}} \\
    \midrule
    Position & 0.189 & 0.176 & 0.114 \\
    Verbosité & 0.029 & 0.012 & 0.004 \\
    Sentiment & 0.073 & 0.104 & 0.267 \\
    Compassion-fade & 0.060 & 0.068 & -0.009 \\
    Bandwagon & 0.100 & 0.046 & 0.025 \\
    Autorité & 0.041 & -0.003 & -0.001 \\
    Distraction & 0.011 & 0.003 & 0.040 \\
    Diversité & -0.015 & 0.019 & -0.019 \\
    \bottomrule
  \end{tabular}
  \caption*{\footnotesize Contrairement à l'article, qui utilise un seul CR par famille de données, nous calculons un CR par biais (voir~\cref{sec-bis-method}).}
\end{table}

